% !PRINTED 383
\weberpage{383}
%<<CONT
\gesperrt{a) Jede Permutation $\pi$ der Wurzeln der Doppeltangentengleichung, die auf alle rationalen Gleichungen zwischen diesen Wurzeln anwendbar ist, gehört zu $P$.}

Dies ergiebt sich so: Wenn $\pi$ auf alle rationalen Gleichungen zwischen den Wurzeln anwendbar ist, so gilt dasselbe von den Potenzen von $\pi$. Nach 1. kann niemals durch $\pi$ ein syzygetisches Tripel in ein azygetisches oder umgekehrt übergeführt werden; denn $\pi$ ist, wenn $\xi_1\xi_2\xi_3$ ein syzygetisches Tripel ist, auf die Gleichung (3) anwendbar; also kann $\xi_1\xi_2\xi_3$ nicht in ein azygetisches Tripel übergehen. Und auch das Umgekehrte ist nicht möglich, weil sonst durch $\pi^{-1}$ ein syzygetisches in ein azygetisches Tripel übergeführt würde. Daher geht auch durch $\pi$ irgend ein vollständiges Siebener-System wieder in ein solches System in irgend welcher Anordnung über, und wenn man dann $\pi$ auf alle Gleichungen von der Form (5) anwendet, so ergiebt sich eine Permutation der $\xi_i, \xi_{ik}$, die zu $P$ gehört.

\gesperrt{b) Jede rationale Gleichung zwischen den Wurzeln der Doppeltangentengleichung gestattet alle Permutationen der Gruppe $P$.}

Eine rationale Relation zwischen den Wurzeln hat die Form

\weq{6}{\Phi(\xi_1,\dots,\xi_7,\xi_{12},\dots,a\dots) = 0,}

worin $\Phi$ eine rationale Function ist, und $a\dots$ die Verhältnisse der Coëfficienten von $f$ bedeuten. Hierin kann man durch (5) die $\xi_{12},\xi_{13},\dots$ rational durch $\xi_1,\xi_2,\dots,\xi_7$ ausdrücken, und nach \S.\,100, II., nachdem die Curve $f$ durch ein vollständiges Siebener-System ihrer Doppeltangenten rational bestimmt ist, lassen sich dann die $a$ rational (mit nur numerischen Coëfficienten) durch die sieben Grössenpaare (2)

\[
\xi_1, \Theta_1;\ \xi_2, \Theta_2;\ \dots;\ \xi_7, \Theta_7
\]

ausdrücken. Da diese aber ganz beliebig gegeben sein können, so muss die Gleichung (6) durch diese Substitutionen in eine Identität übergehen. Die Gleichung (6) muss also auch richtig bleiben, wenn man darin die $\xi_1,\xi_2,\dots,\xi_7$ durch ein anderes Siebener-System oder auch durch dasselbe in anderer Ordnung ersetzt, und gleichzeitig unter den $\xi_{12},\dots$ die durch die Formel (5) vorgeschriebene Permutation vornimmt, d.\,h. wenn man unter den Wurzeln der Doppeltangentengleichung eine Permutation der Gruppe $P$ ausführt.
